\section{}

\paragraph{1222年}\eventanchor{JIN-1222-REGNAL-YEAR}{JIN-1222-REGNAL-YEAR}　金以金宣宗在位第9年作为诏令、赋役与外交文书的纪年框架。账本所列《金史》〈本纪〉、〈交聘表〉及相关志；《宋史》与《建炎以来系年要录》相关条为本条来源起点；该条可固定编年位置，具体执行、规模和地方社会影响仍须保留来源差异。

\paragraph{1223年}\eventanchor{JIN-1223-REGNAL-YEAR}{JIN-1223-REGNAL-YEAR}　金以金宣宗在位第10年作为诏令、赋役与外交文书的纪年框架。账本所列《金史》〈本纪〉、〈交聘表〉及相关志；《宋史》与《建炎以来系年要录》相关条为本条来源起点；该条可固定编年位置，具体执行、规模和地方社会影响仍须保留来源差异。

\paragraph{1224年}\eventanchor{JIN-1224-EVENT-106}{JIN-1224-EVENT-106}　金哀宗即位。账本所列《金史》〈本纪〉、〈交聘表〉及相关志；《宋史》与《建炎以来系年要录》相关条为本条来源起点；该条可固定编年位置，具体执行、规模和地方社会影响仍须保留来源差异。

\paragraph{1224年}\eventanchor{JIN-1224-REGNAL-YEAR}{JIN-1224-REGNAL-YEAR}　金以金宣宗在位第11年作为诏令、赋役与外交文书的纪年框架。账本所列《金史》〈本纪〉、〈交聘表〉及相关志；《宋史》与《建炎以来系年要录》相关条为本条来源起点；该条可固定编年位置，具体执行、规模和地方社会影响仍须保留来源差异。

\paragraph{1225年}\eventanchor{JIN-1225-REGNAL-YEAR}{JIN-1225-REGNAL-YEAR}　金以金哀宗在位第1年作为诏令、赋役与外交文书的纪年框架。账本所列《金史》〈本纪〉、〈交聘表〉及相关志；《宋史》与《建炎以来系年要录》相关条为本条来源起点；该条可固定编年位置，具体执行、规模和地方社会影响仍须保留来源差异。

\paragraph{1226年}\eventanchor{JIN-1226-REGNAL-YEAR}{JIN-1226-REGNAL-YEAR}　金以金哀宗在位第2年作为诏令、赋役与外交文书的纪年框架。账本所列《金史》〈本纪〉、〈交聘表〉及相关志；《宋史》与《建炎以来系年要录》相关条为本条来源起点；该条可固定编年位置，具体执行、规模和地方社会影响仍须保留来源差异。

\paragraph{1227年}\eventanchor{JIN-1227-REGNAL-YEAR}{JIN-1227-REGNAL-YEAR}　金以金哀宗在位第3年作为诏令、赋役与外交文书的纪年框架。账本所列《金史》〈本纪〉、〈交聘表〉及相关志；《宋史》与《建炎以来系年要录》相关条为本条来源起点；该条可固定编年位置，具体执行、规模和地方社会影响仍须保留来源差异。

\paragraph{1228年}\eventanchor{JIN-1228-REGNAL-YEAR}{JIN-1228-REGNAL-YEAR}　金以金哀宗在位第4年作为诏令、赋役与外交文书的纪年框架。账本所列《金史》〈本纪〉、〈交聘表〉及相关志；《宋史》与《建炎以来系年要录》相关条为本条来源起点；该条可固定编年位置，具体执行、规模和地方社会影响仍须保留来源差异。

\paragraph{1229年}\eventanchor{JIN-1229-REGNAL-YEAR}{JIN-1229-REGNAL-YEAR}　金以金哀宗在位第5年作为诏令、赋役与外交文书的纪年框架。账本所列《金史》〈本纪〉、〈交聘表〉及相关志；《宋史》与《建炎以来系年要录》相关条为本条来源起点；该条可固定编年位置，具体执行、规模和地方社会影响仍须保留来源差异。

\paragraph{1230年}\eventanchor{JIN-1230-EVENT-107}{JIN-1230-EVENT-107}　蒙古重新集中兵力攻金。账本所列《金史》〈本纪〉、〈交聘表〉及相关志；《宋史》与《建炎以来系年要录》相关条为本条来源起点；该条可固定编年位置，具体执行、规模和地方社会影响仍须保留来源差异。

\paragraph{1230年}\eventanchor{JIN-1230-REGNAL-YEAR}{JIN-1230-REGNAL-YEAR}　金以金哀宗在位第6年作为诏令、赋役与外交文书的纪年框架。账本所列《金史》〈本纪〉、〈交聘表〉及相关志；《宋史》与《建炎以来系年要录》相关条为本条来源起点；该条可固定编年位置，具体执行、规模和地方社会影响仍须保留来源差异。

\paragraph{1231年}\eventanchor{JIN-1231-REGNAL-YEAR}{JIN-1231-REGNAL-YEAR}　金以金哀宗在位第7年作为诏令、赋役与外交文书的纪年框架。账本所列《金史》〈本纪〉、〈交聘表〉及相关志；《宋史》与《建炎以来系年要录》相关条为本条来源起点；该条可固定编年位置，具体执行、规模和地方社会影响仍须保留来源差异。

\paragraph{1232年}\eventanchor{JIN-1232-EVENT-108}{JIN-1232-EVENT-108}　蒙古围攻汴京。账本所列《金史》〈本纪〉、〈交聘表〉及相关志；《宋史》与《建炎以来系年要录》相关条为本条来源起点；该条可固定编年位置，具体执行、规模和地方社会影响仍须保留来源差异。

\paragraph{1232年}\eventanchor{JIN-1232-REGNAL-YEAR}{JIN-1232-REGNAL-YEAR}　金以金哀宗在位第8年作为诏令、赋役与外交文书的纪年框架。账本所列《金史》〈本纪〉、〈交聘表〉及相关志；《宋史》与《建炎以来系年要录》相关条为本条来源起点；该条可固定编年位置，具体执行、规模和地方社会影响仍须保留来源差异。

\paragraph{1233年}\eventanchor{JIN-1233-EVENT-109}{JIN-1233-EVENT-109}　金哀宗退至蔡州。账本所列《金史》〈本纪〉、〈交聘表〉及相关志；《宋史》与《建炎以来系年要录》相关条为本条来源起点；该条可固定编年位置，具体执行、规模和地方社会影响仍须保留来源差异。

\paragraph{1233年}\eventanchor{JIN-1233-REGNAL-YEAR}{JIN-1233-REGNAL-YEAR}　金以金哀宗在位第9年作为诏令、赋役与外交文书的纪年框架。账本所列《金史》〈本纪〉、〈交聘表〉及相关志；《宋史》与《建炎以来系年要录》相关条为本条来源起点；该条可固定编年位置，具体执行、规模和地方社会影响仍须保留来源差异。

\paragraph{1234年}\eventanchor{JIN-1234-EVENT-110}{JIN-1234-EVENT-110}　蒙古与南宋军攻破蔡州，金朝灭亡。账本所列《金史》〈本纪〉、〈交聘表〉及相关志；《宋史》与《建炎以来系年要录》相关条为本条来源起点；该条可固定编年位置，具体执行、规模和地方社会影响仍须保留来源差异。

\paragraph{1234年}\eventanchor{JIN-1234-REGNAL-YEAR}{JIN-1234-REGNAL-YEAR}　金以金哀宗在位第10年作为诏令、赋役与外交文书的纪年框架。账本所列《金史》〈本纪〉、〈交聘表〉及相关志；《宋史》与《建炎以来系年要录》相关条为本条来源起点；该条可固定编年位置，具体执行、规模和地方社会影响仍须保留来源差异。

